\subsection{A formulation-specific literature comparison}
Table~\ref{tab:literature} distinguishes constructions, not architecture names alone. A stated constraint is not a universal claim about every network bearing the same acronym. Conversely, the absence of a constraint does not establish failure on every deformation path. The accompanying audit records the consulted source versions and separates complete text reading from inspection of the specific passages supporting a claim.

\begingroup
\footnotesize
\setlength{\tabcolsep}{4pt}
\renewcommand{\arraystretch}{1.12}
\begin{longtable}{@{}>{\raggedright\arraybackslash}p{2.6cm}>{\raggedright\arraybackslash}p{4.05cm}>{\raggedright\arraybackslash}p{4.15cm}>{\raggedright\arraybackslash}p{4.85cm}@{}}
\caption{Constitutive and reduced-model antecedents: variables, imposed structure, and limits of the comparison. Statements refer to the specified formulation and its hypotheses.}
\label{tab:literature}\\
\toprule
Source / variant & Represented object and variables & Relevant imposed structure & Distinction from a broader guarantee\\
\midrule
\endfirsthead
\multicolumn{4}{l}{Table \thetable\ continued}\\
\toprule
Source / variant & Represented object and variables & Relevant imposed structure & Distinction from a broader guarantee\\
\midrule
\endhead
\bottomrule
\endfoot
As'ad et al.~\cite{asad2022}
& Scalar energy in objective Green strain; stress from differentiation
& Potential structure; ICNN convexity in Green strain
& Stability argument depends on stress state. Convexity in strain is not a global polyconvexity certificate.\\
\addlinespace
Xu et al.~\cite{Xu2021}
& Incremental stress update using a learned lower-triangular matrix factor
& The product $\bm L\bm L^T$ is positive semidefinite, and positive definite if $\bm L$ is nonsingular
& The factorized matrix is not necessarily an energy Hessian. Differentiating a state-dependent incremental update includes additional terms.\\
\addlinespace
Klein et al.~\cite{klein2022}: invariant model
& Convex nondecreasing energy network acting on polyconvex invariants
& Objectivity and polyconvexity by construction for the selected invariants
& An admissible invariant set may restrict approximation. Its completeness does not follow from core-network capacity.\\
\addlinespace
Klein et al.~\cite{klein2022}: gradient model
& Network of deformation minors with material-group symmetrization
& Polyconvexity and exact selected finite-group symmetry
& Objectivity is learned by augmentation. Continuous transverse isotropy is approximated by a finite subgroup in the reported example.\\
\addlinespace
Linden et al.~\cite{linden2023}
& Invariant energy network with analytic growth and reference corrections
& Compatible polyconvex normalization for the stated isotropic and transverse-isotropic cases
& Global nonnegative energy is examined separately. Zero energy and stress at the identity alone do not prove it.\\
\addlinespace
Linka and Kuhl~\cite{Linka2023}
& CANN using invariant energy terms and interpretable learned coefficients
& Architecture encodes constitutive restrictions for the specified building blocks
& Automated selection within a candidate family is distinct from completeness of that family for arbitrary anisotropy.\\
\addlinespace
Tac et al.~\cite{Tac2022}
& Scalar NODE flows construct monotone invariant-energy derivatives
& Integrable invariant expansion and polyconvexity under the stated nonnegative/monotone construction
& Worked tissue examples use selected fiber invariants; the framework also discusses extensions. It should not be called incapable of compressible modeling.\\
\addlinespace
Ghaderi et al.~\cite{Ghaderi2020}
& Directional learning agents, microsphere summation, and maximum-history variables
& Physics-based decomposition for polymer elasticity and inelastic effects
& This is not a purely history-independent constitutive tier. The present comparison does not adopt a global convexity theorem merely from signs of network weights.\\
\addlinespace
Abdolazizi et al.~\cite{Abdolazizi2025}
& CKAN scalar energy; spline edges, reference correction, and symbolic simplification
& Energy differentiation and partial monotonicity in selected inputs
& Partial monotonicity is not input convexity. Section 6 identifies polyconvexity as a further extension; CKAN and ICKAN are not equivalent.\\
\addlinespace
Kalina et al.~\cite{Kalina2024}
& Invariants formed with learned generalized structure tensors of orders two, four, or six
& Objective energy representation and identification of anisotropy through the tensor construction
& Learning anisotropy is an existing capability. A general structure-tensor invariant is not automatically a polyconvex feature.\\
\addlinespace
Klein et al.~\cite{klein2026}
& Objective triclinic polyconvex features, with optional finite-group symmetrization
& Polyconvexity and objectivity without requiring a prescribed nontrivial material symmetry
& Its formulation is not restricted to the cubic numerical examples. Absence of a prescribed group is not the novelty claimed here.\\
\addlinespace
Thakolkaran et al.~\cite{thakolkaran2025}
& Input-convex monotone spline KAN for isotropic compressible energy
& Polyconvex composition on the constrained spline domains
& Exact endpoint tangency must be distinguished from finite-secant tail formulas. The present integrated-hat spline has analytic $C^2$ joins.\\
\addlinespace
Barnett et al.~\cite{barnett2023}; Ares De Parga et al.~\cite{aresdeparga2026nonlinear}
& Primary/secondary state partition; secondary coordinates learned from primary coordinates
& Low-dimensional closure integrated into a projection-based equilibrium or evolution model
& Learning a closure does not remove the reduced solve. The direct strain substitution used here is a separate operating mode.\\
\addlinespace
Farhat et al.~\cite{farhat2015}
& Fixed-weight ECSW applied to reduced nonlinear finite element dynamics
& Lagrangian structure preservation under the specified hypotheses
& The result does not automatically cover state-dependent weights and an independently approximated macroscopic stress output.\\
\addlinespace
Hern\'andez et al.~\cite{hernandez2026}
& Nonlinear-manifold virtual work with manifold-adaptive empirical cubature
& Adaptive integration constraints and reduced equilibrium
& Weight positivity and integration accuracy are distinct. This paper uses that methodology, not a new ECM algorithm.\\
\end{longtable}
\endgroup

The comparison suggests two axes rather than a single ranking. Constitutive constraints determine the admissible family of stress--strain maps; reduction choices determine which microscopic information and equations are retained online. An HPROM can be more accurate than a polyconvex network without inheriting its macroscopic constitutive certificate. An unrestricted energy network can fit a test set well while lacking a global rank-one-curvature constraint. Both statements are compatible with the results reported here.
